\chapter{The Agent Loop Architecture}
\label{ch:agent-loop}

\epigraph{``The loop is not the implementation. The loop is the contract.''}{}

\learninggoal{
	\item Decompose the agent loop into its essential phases.
	\item Understand the message-passing model that underlies all agent frameworks.
	\item Distinguish between reactive, deliberative, and reflective loop variants.
	\item Compare state management strategies across frameworks.
}

\section{The Universal Loop}

Every agent---from a single-prompt tool caller to a multi-agent swarm---implements the same abstract loop. The implementation differs, but the phases are universal:

\begin{enumerate}
	\item \textbf{Observe:} read the current state of the environment (sensors, API responses, user input).
	\item \textbf{Think:} reason about what action to take given the current state and goal.
	\item \textbf{Act:} execute the chosen action (tool call, response, internal state update).
	\item \textbf{Reflect:} evaluate the outcome and adjust future behaviour.
\end{enumerate}

\begin{definition}[Agent Loop]
	An \textbf{agent loop} is a control structure that repeatedly invokes an \llm\ as a reasoning engine, passing in observations and memory, parsing structured actions from the output, executing them, and incorporating results into the next iteration.
\end{definition}

\section{The Message-Passing Model}

The cleanest abstraction for an agent loop is \emph{message passing}. Every component---LLM, tool, memory store, user---is a node that sends and receives messages.

\begin{lstlisting}[caption=Message types in the agent loop]
  // Core message types
  Message ::= {
    role:    "system" | "user" | "assistant" | "tool",
    content: string,
    tool_calls?: ToolCall[],   // assistant only
    tool_result?: string,       // tool only
    metadata?:  { key: value }
  }

  // Loop invariant
  messages[0].role == "system"  // first message is the system prompt
  messages[-1].role != "assistant"  // last message is pending
\end{lstlisting}

This model has several desirable properties:
\begin{itemize}
	\item \textbf{Immutability:} old messages are never modified, only appended. This enables replay, debugging, and checkpointing.
	\item \textbf{Determinism:} given the same message history and seed, the same actions are produced.
	\item \textbf{Inspectability:} the full decision trace is available for logging and evaluation.
	\item \textbf{Tool integration:} tool results are just messages with a special \texttt{role}.
\end{itemize}

\section{Loop Variants}

\subsection{Reactive Loop}

The simplest variant: observe $\to$ act, no explicit reasoning step.

\begin{lstlisting}[caption=Reactive agent loop]
  messages = [system_prompt]
  while not done:
      user_input = input()
      messages.append({"role": "user", "content": user_input})
      response = llm(messages)
      messages.append({"role": "assistant", "content": response})
\end{lstlisting}

This is effectively a chatbot. It is an agent only if the response can trigger side effects (e.g., ``send this email'').

\subsection{Reason-Act Loop (ReAct)}

The canonical agent loop: interleave reasoning traces with actions.

\begin{lstlisting}[caption=ReAct-style loop]
  messages = [system_prompt, user_query]
  while not done:
      response = llm(messages)           // Think
      action   = parse_action(response)   // Parse structured output
      if action.type == "final_answer":
          return action.content           // Done
      result   = execute_tool(action)     // Act
      messages.append(response)
      messages.append({"role": "tool",   // Observe
                       "content": result})
\end{lstlisting}

This was formalised by Yao et al.~\cite{yao2023react} and is the basis of most production agent systems. The key insight is that the reasoning trace (\emph{thought}) and the action are generated in the same forward pass, eliminating the overhead of a separate planner.

\subsection{Reflective Loop (Reflexion)}

Add an explicit evaluation step after action execution:

\begin{lstlisting}[caption=Reflexion-style loop]
  while not done:
      trajectory = []
      while not subtask_done:
          action = reason_and_act(state)
          trajectory.append(action)
      reflection = reflect(trajectory)     // Evaluate
      memory.add(reflection)               // Learn
\end{lstlisting}

Reflexion~\cite{shinn2023reflexion} adds a \texttt{reflect} step that generates a natural-language summary of what went wrong and what to change. This summary is stored in memory and influences future episodes. It converts trial-and-error into verbal reinforcement learning.

\subsection{Hierarchical Loop}

For complex tasks, a single loop is insufficient. Hierarchical agents use a manager--worker pattern:

\begin{lstlisting}[caption=Hierarchical agent loops]
  manager_loop():
      goal = receive_goal()
      plan = decompose(goal)            // Break into sub-tasks
      for subtask in plan:
          result = worker_loop(subtask) // Delegate to sub-agent
          verify(result, goal)
      return assemble(plan, results)

  worker_loop(subtask):
      for step in subtask.steps:
          action = reason(step)
          execute(action)
      return output
\end{lstlisting}

The key design decision is the \emph{granularity of decomposition}: too coarse and workers are overloaded; too fine and coordination overhead dominates.

\section{State Management}

Every agent loop must maintain state across iterations. The state includes:

\begin{itemize}
	\item \textbf{Message history:} the full conversation (grows unboundedly).
	\item \textbf{Execution context:} variables, file handles, authentication tokens.
	\item \textbf{Plan state:} which steps are done, which remain.
	\item \textbf{Memory:} long-term knowledge stored beyond the current session.
\end{itemize}

\subsection{Context Window Management}

The message history grows with every loop iteration, eventually exceeding the \llm 's context window. Common strategies:

\begin{longtable}{p{3.5cm}p{4cm}p{4cm}p{3cm}}
	\toprule
	\textbf{Strategy}   & \textbf{Mechanism}                  & \textbf{Trade-off}         & \textbf{Used by}  \\
	\midrule
	Sliding window      & Drop oldest messages                & Loses early context        & Simple chatbots   \\
	Summary compression & Replace old messages with a summary & Loses detail, adds latency & LangChain         \\
	Token budget        & Evict least-relevant messages       & Relevance heuristic needed & Custom            \\
	Structured memory   & Store key info in external DB       & Requires retrieval infra   & Production agents \\
	\bottomrule
\end{longtable}

\begin{principle}[Token Budgeting]
	An agent should never exceed 80\% of its context window during normal operation. The remaining 20\% is a safety margin for tool outputs and unexpected long responses.
\end{principle}

\subsection{Checkpointing}

Production agents checkpoint their state at every loop iteration:

\begin{lstlisting}[caption=Agent checkpoint structure]
  Checkpoint ::= {
    id:            UUID,
    messages:      Message[],
    memory_snapshot: { key: value },
    plan_state:    { completed: Step[], pending: Step[] },
    tool_states:   { connection_id: state },
    created_at:    timestamp
  }
\end{lstlisting}

Checkpoints enable:
\begin{itemize}
	\item \textbf{Fault recovery:} restart from the last checkpoint after a crash.
	\item \textbf{Forking:} try alternative actions from a given point.
	\item \textbf{Audit:} replay the full decision trace.
\end{itemize}

\section{Framework Comparison}

All agent frameworks implement the same core loop but differ in how they structure reasoning, tools, and memory.

\begin{longtable}{p{2.5cm}p{3.5cm}p{3.5cm}p{3.5cm}}
	\toprule
	\textbf{Framework} & \textbf{Loop model}          & \textbf{Tool model}                 & \textbf{State management}           \\
	\midrule
	LangChain          & ReAct-style graph            & Decorated functions, structured I/O & Message window, optional checkpoint \\
	AutoGen            & Multi-agent conversation     & Agent-registered functions          & Full message history                \\
	CrewAI             & Hierarchical manager--worker & Per-role tool sets                  & MongoDB checkpointing               \\
	Semantic Kernel    & Stepwise planner             & Plugin registry, OpenAPI import     & Volatile by default                 \\
	OpenAI Assistants  & Managed loop (hosted)        & Function definition API             & Server-side checkpointing           \\
	\bottomrule
\end{longtable}

\section{Exercises}

\begin{exercise}
	Implement a ReAct agent loop from scratch in pseudocode (no framework). Your loop must support: multiple tools, error handling for tool failures, and a maximum iteration limit.
\end{exercise}

\begin{exercise}
	Compare the context-window management strategies above for an agent that processes a 10-turn conversation with 3 tool calls per turn. Compute total token usage for each strategy.
\end{exercise}

\begin{exercise}
	Design a checkpoint format for an agent that manages a multi-file code change. What state must be saved beyond the message history?
\end{exercise}
